\section{INTRODUCCIÓN}

El avance acelerado de la computación cuántica constituye una de las amenazas más significativas para la infraestructura de seguridad digital contemporánea. Los sistemas de votación electrónica que actualmente se emplean en distintos procesos institucionales dependen, en su mayoría, de criptosistemas asimétricos convencionales como RSA y ECDSA, cuya seguridad reside en la dificultad computacional de factorizar números enteros grandes o de calcular logaritmos discretos sobre curvas elípticas. Esta dependencia, sostenida durante más de tres décadas, se ve hoy cuestionada por la evolución de la computación cuántica, que amenaza con volver obsoletas las garantías criptográficas sobre las cuales se ha edificado la confianza digital moderna.\

La implementación del algoritmo de Shor en computadoras cuánticas de escala relevante permitirá, en un horizonte temporal aún incierto pero cada vez más próximo, romper los esquemas de clave pública convencionales en tiempo polinomial, lo cual invalida las garantías de integridad y secreto del sufragio en los sistemas de votación electrónica vigentes. Este escenario ha sido reconocido de manera formal por organismos de estandarización internacional, entre ellos el National Institute of Standards and Technology (NIST), que desde 2016 conduce un proceso sistemático de evaluación y selección de algoritmos resistentes a ataques cuánticos, cuyo estado de avance ha sido documentado en sucesivos informes técnicos (Alagic et al., 2022). Frente a este panorama, resulta imperativa la transición hacia la criptografía post-cuántica (PQC) en todos los dominios de aplicación crítica, y de manera particular en aquellos vinculados a procesos democráticos, donde la confianza pública en la integridad del sufragio constituye un pilar de la estabilidad institucional.\

La urgencia de esta transición no radica únicamente en la disponibilidad futura de computadoras cuánticas suficientemente potentes, sino también en la existencia de ataques de tipo "cosechar ahora, descifrar después" (Store-Now-Decrypt-Later), mediante los cuales adversarios con capacidad de almacenamiento pueden interceptar y conservar votos cifrados en la actualidad, con la expectativa de descifrarlos en el futuro una vez que dispongan de hardware cuántico apropiado. Este riesgo ha sido señalado como crítico para los sectores financiero y gubernamental por organismos como el FS-ISAC (2023) y ha motivado hojas de ruta coordinadas de transición en distintas regiones, incluyendo la Unión Europea (Encryption Consulting, 2026) y el bloque del G7 (G7 Cyber Expert Group, 2026).\

En este contexto, la presente investigación propone el diseño e implementación de un protocolo de votación electrónica con verificabilidad de extremo a extremo (End-to-End Verifiability, E2E) que integra algoritmos criptográficos basados en retículos matemáticos (lattices), recientemente estandarizados por el NIST bajo los marcos FIPS 203 y FIPS 204 (National Institute of Standards and Technology \[NIST\], 2024). La solución propuesta se fundamenta en tres pilares complementarios: primero, el uso de ML-KEM para el cifrado del voto y de ML-DSA para la autenticación y firma digital del votante; segundo, la aplicación de pruebas de conocimiento cero de tipo zk-SNARK para preservar la privacidad del elector sin comprometer la verificabilidad pública del proceso; y tercero, la utilización de una infraestructura blockchain que actúa como tablero de anuncios (Bulletin Board) inmutable y de acceso público, sobre el cual se ejecutan contratos inteligentes encargados de validar de forma autónoma tanto las pruebas criptográficas como las firmas post-cuánticas.\

A diferencia de propuestas puramente teóricas que han dominado la literatura reciente sobre votación post-cuántica, el sistema que aquí se plantea está explícitamente optimizado para ejecutarse sobre hardware convencional de arquitectura x86 o ARM, sin requerir dispositivos criptográficos especializados ni infraestructura de cómputo de alto rendimiento. Los resultados preliminares de viabilidad técnica, obtenidos durante la fase exploratoria del proyecto, muestran tiempos de procesamiento del orden de 164.9 milisegundos por voto, cifra que resulta consistente con los hallazgos reportados por estudios recientes sobre la evaluación de algoritmos post-cuánticos en dispositivos de recursos limitados (Lopez, Cadena y Rahman, 2025), y que demuestra que la seguridad cuántica no impone necesariamente una penalización prohibitiva en términos de rendimiento.\

El presente documento constituye el perfil de tesis de grado que fundamenta esta investigación. En él se exponen los antecedentes que enmarcan el problema, se define el planteamiento del problema central y los problemas secundarios asociados, se establecen los objetivos general y específicos, se formula la hipótesis de trabajo junto con la operacionalización de sus variables, se justifica la pertinencia económica, social y científica del estudio, se delimitan los alcances y límites de la propuesta, se detallan los aportes prácticos y teóricos esperados, se describe la metodología a emplear, se desarrolla el marco teórico que sustenta las decisiones de diseño, y finalmente se presenta el índice tentativo, el cronograma de avance y la bibliografía consultada.\

